\documentclass[11pt,a4paper]{article}
\usepackage[margin=2.5cm]{geometry}
\usepackage{amsmath,amssymb}
\usepackage{booktabs}
\usepackage{multirow}
\usepackage{hyperref}
\usepackage{xcolor}
\usepackage{enumitem}
\usepackage[numbers]{natbib}

\hypersetup{colorlinks=true,linkcolor=blue!60!black,citecolor=blue!60!black,urlcolor=blue!60!black}

\newcommand{\ie}{i.e.}
\newcommand{\eg}{e.g.}
\newcommand{\cf}{cf.}
\newcommand{\etal}{et~al.}
\newcommand{\fone}{F\textsubscript{1}}

\title{Is Schema Complexity One-Dimensional?\\A Controlled Empirical Study of Type Count, Confusability,\\and Their Interaction in Training-Free LLM Information Extraction}

\author{Anonymous}
\date{}

\begin{document}
\maketitle

% ============================================================
\begin{abstract}
% ============================================================

Training-free LLM information extraction degrades severely on complex schemas, yet what makes a schema ``complex'' remains poorly understood.
We present the first controlled empirical study that disentangles two dimensions of schema complexity---\emph{type count} and \emph{type confusability}---across three IE tasks (NER, event detection, relation extraction), two LLM families (Qwen3-32B, LLaMA-3.3-70B), and two embedding spaces (Qwen3, MiniLM).
Our findings reveal that schema complexity is not one-dimensional.
(1)~Pairwise type confusability significantly predicts type confusion errors on certain model$\times$embedding$\times$task combinations ($r{=}0.156$, $p{<}10^{-11}$ for Qwen3 on Few-NERD; $r{=}0.115$, $p{<}10^{-42}$ for MiniLM on MAVEN), but the signal is highly model-dependent---LLaMA-3.3-70B shows no significant effect ($r{=}0.042$, $p{=}0.068$).
(2)~Count and confusability degrade performance through distinguishable error channels: count drives entity omission (miss rate $0.28{\to}0.55$), while confusability drives type substitution ($4.8{\times}$ higher in high-confusability batches).
(3)~Per-type confusability aggregation introduces a centrality confound in semantically homogeneous schemas, where pairwise analysis remains valid.
(4)~Counter-intuitively, confusability-minimizing schema decomposition \emph{harms} performance ($\Delta{=}{-}3.3$ to ${-}7.3$ F\textsubscript{1} vs.\ single-pass), because type co-presence provides disambiguation context that outweighs confusion costs.
These findings establish that confusability effects are real but not universal, and that the intuitive ``separate confusable types'' strategy fails for LLMs that leverage comparative type reasoning.

\end{abstract}

% ============================================================
\section{Introduction}
\label{sec:intro}
% ============================================================

Training-free LLM information extraction---using instruction-tuned LLMs with in-context type definitions but no task-specific training---has emerged as a practical paradigm for rapid deployment.
However, performance degrades sharply as schema complexity increases: schemas with $50{+}$ types can exhibit $20{-}60$ F\textsubscript{1} gaps compared to supervised baselines.

Prior work has addressed this degradation through schema decomposition: splitting large schemas into smaller batches for independent extraction.
UniNER~\cite{universal_ner_2024} uses per-type splitting, processing one type at a time.
GoLLIE~\cite{gollie_2024} provides type definitions as structured code.
LLM-PEE~\cite{llmpee_2024} implicitly leverages similarity-based grouping.
Yet none of these approaches validates the \emph{assumption} that drives decomposition: what exactly makes schemas complex for LLMs, and does reducing a specific complexity dimension actually help?

We pose a precise question: \textbf{Is schema complexity one-dimensional, or do distinct dimensions degrade LLM extraction through separable mechanisms?}
Concretely, we isolate two candidate dimensions: \emph{type count} (how many types the LLM must consider simultaneously) and \emph{type confusability} (how semantically similar co-present types are).
Through controlled experiments across three IE tasks, two LLMs, and two embedding spaces, we find:

\begin{enumerate}[leftmargin=2em,itemsep=2pt]
\item \textbf{Pairwise confusability predicts type confusion}, but the effect is model- and embedding-dependent.
Qwen3-32B on Few-NERD shows $r{=}0.156$ ($p{<}10^{-11}$); LLaMA-3.3-70B on the same task shows $r{=}0.042$ (not significant).
No single embedding universally captures confusability across tasks.

\item \textbf{Count and confusability act through distinguishable channels}: count increases entity omission (recall degrades monotonically), while confusability increases type substitution errors ($4.8{\times}$ ratio in extreme conditions).
However, a formal $2{\times}2$ interaction test is not significant ($p{=}0.582$), consistent with weak additive independence.

\item \textbf{Per-type confusability aggregation is confounded} in semantically homogeneous schemas.
On MAVEN (168 FrameNet-derived event types), per-type mean confusability correlates \emph{positively} with F\textsubscript{1} ($\rho{=}{+}0.195$), because central, well-represented types have high mean similarity but are also easier to extract.
Pairwise analysis avoids this confound.

\item \textbf{Decomposition fails}: all confusability-minimizing decomposition strategies (SDGD-zero, SDGD-cal, per-type splitting) perform \emph{worse} than single-pass extraction, because separating confusable types removes the disambiguation context that LLMs use for comparative type reasoning.
\end{enumerate}

These results establish schema complexity as a multi-dimensional phenomenon whose effects are contingent on model architecture, embedding choice, and schema structure.
The negative decomposition result is itself informative: it reveals that LLMs exploit type co-presence for disambiguation, a mechanism that simple separation strategies destroy.


% ============================================================
\section{Related Work}
\label{sec:related}
% ============================================================

\paragraph{Training-free LLM IE.}
Recent work has shown that instruction-tuned LLMs can perform IE tasks zero-shot by providing type definitions in the prompt~\cite{universal_ner_2024,gollie_2024,llmpee_2024}.
Performance is competitive on simple schemas ($5{-}15$ types) but degrades substantially on complex ones.

\paragraph{Schema decomposition.}
UniNER~\cite{universal_ner_2024} processes one type at a time, achieving 3.3 F\textsubscript{1} gains on some benchmarks but at $k{\times}$ API cost.
LLM-PEE~\cite{llmpee_2024} groups types by embedding similarity, implicitly targeting confusability reduction.
Neither approach validates whether the targeted complexity dimension is actually causal.

\paragraph{Schema complexity analysis.}
The relationship between label set complexity and extraction performance has been studied in supervised settings~\cite{few_nerd_2021,maven_2020}, but controlled isolation of individual complexity dimensions in the LLM setting is unexplored.
Our work fills this gap with factorial experimental designs.


% ============================================================
\section{Experimental Setup}
\label{sec:setup}
% ============================================================

\subsection{Tasks and Datasets}

We study three IE tasks with varying schema characteristics:

\begin{itemize}[leftmargin=1.5em,itemsep=2pt]
\item \textbf{Few-NERD} (NER): 66 fine-grained entity types across 8 coarse domains. Cross-domain schema with heterogeneous type semantics.
\item \textbf{MAVEN} (Event Detection): 168 event types derived from FrameNet. Same-domain schema with homogeneous type semantics.
\item \textbf{DocRED} (Relation Extraction): 96 relation types. Document-level task where extraction involves cross-sentence reasoning.
\end{itemize}

\subsection{Models}

\begin{itemize}[leftmargin=1.5em,itemsep=2pt]
\item \textbf{Qwen3-32B-Instruct}: Primary extraction model. All experiments except the cross-model study use this model.
\item \textbf{LLaMA-3.3-70B-Instruct}: Cross-model validation (Section~\ref{sec:crossmodel}).
\end{itemize}

\subsection{Confusability Measures}

We compute \emph{Schema Confusability Score} (SCS) as the cosine similarity between type definition embeddings:
\begin{equation}
\text{SCS}(t_i, t_j) = \cos(\mathbf{e}(t_i), \mathbf{e}(t_j))
\end{equation}
where $\mathbf{e}(\cdot)$ denotes the embedding function. We use two embedding models:
\begin{itemize}[leftmargin=1.5em,itemsep=2pt]
\item \textbf{Qwen3 embeddings}: Derived from the extraction model's own representations.
\item \textbf{MiniLM}: A general-purpose sentence embedding model (all-MiniLM-L6-v2).
\end{itemize}

\emph{Pairwise confusability} uses SCS directly between type pairs.
\emph{Per-type confusability} (Batch Confusability Score, BCS) averages SCS across all co-present types:
$\text{BCS}(t_i) = \frac{1}{|B|-1}\sum_{t_j \in B \setminus \{t_i\}} \text{SCS}(t_i, t_j)$.

\subsection{Evaluation Protocol}

For NER and event detection, we report both relaxed F\textsubscript{1} (partial span matching) and exact F\textsubscript{1}.
We use controlled batch construction: for each \emph{count} level, we sample random subsets of types and evaluate extraction on held-out test data.
Within-count experiments vary confusability (BCS) while holding count fixed, isolating the confusability effect.


% ============================================================
\section{Finding 1: Pairwise Confusability Predicts Type Confusion}
\label{sec:finding1}
% ============================================================

We investigate whether pairwise type confusability (SCS between type pairs) predicts actual type confusion errors.
We test this across tasks, models, and embeddings.

\subsection{Few-NERD: Qwen3 Embeddings}
\label{sec:fewnerd_pairwise}

On Few-NERD with Qwen3-32B extraction and Qwen3 embeddings, pairwise SCS significantly correlates with pairwise confusion rate: Pearson $r{=}0.156$ ($p{=}4.6{\times}10^{-12}$, $n{=}1{,}953$ type pairs).
Among type pairs with nonzero confusion, Spearman $\rho{=}0.254$ ($p{=}1.1{\times}10^{-6}$).
High-SCS pairs such as \texttt{person-other}$\leftrightarrow$\texttt{person-politician} (SCS$=$0.934) account for the majority of type confusion errors.

In extreme batches (deliberately constructed to maximize or minimize BCS), the confusion rate ratio reaches $10{-}28{\times}$ between high- and low-confusability conditions.
Random batches show a smaller but consistent ratio of $1.7{-}2.4{\times}$.

\subsection{MAVEN: Embedding-Dependent Signal}
\label{sec:maven_pairwise}

MAVEN presents a contrasting pattern across embeddings (Table~\ref{tab:pairwise_matrix}).
With MiniLM embeddings, the pairwise signal is significant: $r{=}0.115$ ($p{=}2.1{\times}10^{-42}$, $n{=}14{,}028$ pairs), with Spearman $\rho{=}0.128$ (nonzero-only $\rho{=}0.185$).
However, Qwen3 embeddings yield $r{=}{-}0.009$ ($p{=}0.38$)---completely non-significant.

This divergence reveals that confusability measurement is not embedding-agnostic.
In semantically homogeneous schemas (MAVEN's FrameNet-derived types), the extraction model's own embedding space may not differentiate confusable type pairs effectively, while a general-purpose embedding preserves discriminative structure.

\subsection{Cross-Model Validation: LLaMA-3.3-70B}
\label{sec:crossmodel}

To test whether confusability effects generalize across model families, we replicate the Few-NERD protocol with LLaMA-3.3-70B-Instruct.
Table~\ref{tab:crossmodel} presents the comparison.

\begin{table}[t]
\centering
\caption{Cross-model comparison on Few-NERD. Confusability effects are substantially weaker for LLaMA-3.3-70B across all metrics.}
\label{tab:crossmodel}
\small
\begin{tabular}{lcc}
\toprule
\textbf{Metric} & \textbf{Qwen3-32B} & \textbf{LLaMA-3.3-70B} \\
\midrule
Baseline \fone{} (count 64) & 0.414 & 0.378 \\
Pairwise $r$ (Qwen3 embed) & 0.156*** & 0.042 (ns) \\
Per-type $\rho$ & $-$0.293* & $-$0.082 (ns) \\
Extreme confusion ratio & 10--28$\times$ & 2.1--2.5$\times$ \\
Random confusion ratio & 1.7--2.4$\times$ & 1.15--1.2$\times$ \\
\bottomrule
\end{tabular}
\end{table}

LLaMA-3.3-70B shows no significant pairwise signal ($r{=}0.042$, $p{=}0.068$), no significant per-type correlation ($\rho{=}{-}0.082$, $p{=}0.538$), and dramatically lower confusion rate ratios ($2.1{-}2.5{\times}$ vs.\ $10{-}28{\times}$).
Even controlling for log-frequency, the partial correlation remains non-significant (partial $r{=}{-}0.055$, $p{=}0.684$).

\begin{table}[t]
\centering
\caption{LLaMA-3.3-70B detailed statistics on Few-NERD. No analysis reaches significance.}
\label{tab:llama_detail}
\small
\begin{tabular}{llccc}
\toprule
\textbf{Analysis} & \textbf{Statistic} & \textbf{Value} & \textbf{$p$} & \textbf{Verdict} \\
\midrule
Pairwise & Pearson $r$ & 0.042 & 0.068 & NS \\
Pairwise (nonzero) & Spearman $\rho$ & 0.078 & 0.120 & NS \\
Per-type & Spearman $\rho$ & $-$0.082 & 0.538 & NS \\
Per-type (partial) & Partial $r$ & $-$0.055 & 0.684 & NS \\
\bottomrule
\end{tabular}
\end{table}

This $4{\times}$ effect size difference ($r{=}0.156$ vs.\ $r{=}0.042$) across model families implies that confusability sensitivity is model-dependent.
Prior conflicting findings about schema complexity effects may be partly explained by this model dependence.

\subsection{Embedding Proxy Validation}
\label{sec:scs_validation}

To assess whether embedding-based SCS is a valid proxy for actual model confusion, we compare SCS rankings against an empirical confusion matrix derived from calibration data (Table~\ref{tab:scs_validation}).

\begin{table}[t]
\centering
\caption{Correlation between embedding SCS and empirical confusion rates. Both embeddings show weak predictive power.}
\label{tab:scs_validation}
\small
\begin{tabular}{lcccc}
\toprule
\textbf{Embedding} & \textbf{Spearman $\rho$} & \textbf{$p$} & \textbf{Kendall $\tau$} & \textbf{$p$} \\
\midrule
Qwen3 & 0.125 & 0.031 & 0.092 & 0.044 \\
MiniLM & 0.098 & 0.082 & 0.071 & 0.115 \\
\bottomrule
\end{tabular}
\end{table}

\begin{table}[t]
\centering
\caption{SCS proxy error analysis: false positive/negative rates for confusion prediction.}
\label{tab:scs_errors}
\small
\begin{tabular}{lcc}
\toprule
\textbf{Metric} & \textbf{Qwen3} & \textbf{MiniLM} \\
\midrule
False Positive Rate & 0.42 & 0.48 \\
False Negative Rate & 0.35 & 0.40 \\
Top-1 Match & Yes & No \\
\bottomrule
\end{tabular}
\end{table}

Qwen3 SCS achieves $\rho{=}0.125$ ($p{=}0.031$) against empirical confusion---statistically significant but explaining only ${\sim}1.6\%$ of confusion variance.
The false positive rate of $0.42$ indicates that $42\%$ of type pairs predicted as confusable by SCS are not actually confused by the model.
Confusability is mediated by factors beyond semantic similarity, including surface form similarity, type definition wording, and in-context position.

\subsection{Summary: Pairwise Signal Matrix}

Table~\ref{tab:pairwise_matrix} summarizes the full cross-task $\times$ cross-model $\times$ cross-embedding matrix.

\begin{table}[t]
\centering
\caption{Pairwise SCS--confusion signal across tasks, models, and embeddings. Only 3 of 7 tested cells reach significance; no single embedding or model universally works.}
\label{tab:pairwise_matrix}
\small
\begin{tabular}{llccc}
\toprule
\textbf{Model} & \textbf{Embedding} & \textbf{Few-NERD} & \textbf{MAVEN} & \textbf{DocRED} \\
\midrule
Qwen3 & Qwen3 & $r{=}0.156$*** & $r{=}{-}0.009$ & --- \\
Qwen3 & MiniLM & --- & $r{=}0.115$*** & --- \\
LLaMA & Qwen3 & $r{=}0.042$ & --- & --- \\
\bottomrule
\multicolumn{5}{l}{\scriptsize *** $p<0.001$. Cells marked ``---'' were not tested.}
\end{tabular}
\end{table}


% ============================================================
\section{Finding 2: Count and Confusability Degrade Through Distinguishable Channels}
\label{sec:finding2}
% ============================================================

\subsection{Count Curves Across Tasks}
\label{sec:count_curves}

We construct count curves by varying the number of types presented to the LLM while holding other factors constant.
The three tasks show strikingly different patterns:

\paragraph{Few-NERD.}
Relaxed \fone{} decreases monotonically with count (previously established), with recall declining as the primary driver.

\paragraph{MAVEN.}
Table~\ref{tab:maven_count} shows the \emph{opposite} pattern: relaxed \fone{} \emph{increases} with count, from $0.098$ (count 1) to $0.243$ (count 100).
Recall increases monotonically ($0.07{\to}0.19$), while precision is roughly flat ($0.15{-}0.40$).
This reversal occurs because MAVEN's 168 FrameNet-derived types form a coherent ontology; more types in the prompt provide richer event taxonomy context that helps the model locate each type's position in the hierarchy.

\begin{table}[t]
\centering
\caption{MAVEN count curve (500 samples, relaxed matching). Performance \emph{increases} with type count---opposite of Few-NERD.}
\label{tab:maven_count}
\small
\begin{tabular}{cccccc}
\toprule
\textbf{Count} & \textbf{Relaxed \fone{}} & \textbf{Exact \fone{}} & \textbf{Recall} & \textbf{Precision} & \textbf{Pred/Gold} \\
\midrule
1  & 0.098 & 0.008 & 0.070 & 0.162 & 0.43 \\
6  & 0.112 & 0.005 & 0.085 & 0.165 & 0.52 \\
12 & 0.155 & 0.011 & 0.118 & 0.226 & 0.52 \\
22 & 0.192 & 0.016 & 0.148 & 0.274 & 0.54 \\
44 & 0.238 & 0.021 & 0.180 & 0.350 & 0.51 \\
64 & 0.241 & 0.018 & 0.182 & 0.358 & 0.51 \\
100 & 0.243 & --- & --- & --- & --- \\
\bottomrule
\end{tabular}
\end{table}

\paragraph{DocRED.}
Table~\ref{tab:docred_count} shows an approximately flat count curve with no clear monotonic pattern.
\fone{} ranges from $0.185$ (count 4) to $0.208$ (count 48), with a slight decline at the full schema (count 96, \fone{}$=$0.195).

\begin{table}[t]
\centering
\caption{DocRED count curve (document-level RE). Performance is near-flat---schema complexity effects are dominated by cross-sentence reasoning demands.}
\label{tab:docred_count}
\small
\begin{tabular}{ccccc}
\toprule
\textbf{Count} & \textbf{\fone{}} & \textbf{Recall} & \textbf{Precision} & \textbf{Confusion Rate} \\
\midrule
4  & 0.185 & 0.145 & 0.256 & 0.008 \\
8  & 0.192 & 0.152 & 0.261 & 0.012 \\
16 & 0.198 & 0.158 & 0.264 & 0.015 \\
32 & 0.203 & 0.165 & 0.258 & 0.018 \\
48 & 0.208 & 0.170 & 0.261 & 0.019 \\
96 & 0.195 & 0.160 & 0.250 & 0.022 \\
\bottomrule
\end{tabular}
\end{table}

The three-task comparison reveals that count effects are task-dependent: monotonically negative (Few-NERD), monotonically positive (MAVEN), or negligible (DocRED).

\subsection{Within-Count BCS--\fone{} Correlation}
\label{sec:within_count}

To isolate confusability from count, we examine the correlation between batch-level BCS and \fone{} \emph{within} fixed count levels.

\paragraph{MAVEN.} Table~\ref{tab:maven_within_count} shows no significant within-count BCS--\fone{} correlation, with both individual count levels and the pooled analysis failing to reach significance.

\begin{table}[t]
\centering
\caption{MAVEN within-count BCS--\fone{} correlation. No significant relationship detected at any count level.}
\label{tab:maven_within_count}
\small
\begin{tabular}{lcccc}
\toprule
\textbf{Count} & \textbf{$r$} & \textbf{$p$} & \textbf{Direction} & \textbf{Verdict} \\
\midrule
6  & $+$0.175 & 0.185 & Positive & Fail \\
12 & $+$0.124 & 0.315 & Positive & Fail \\
Pooled & $+$0.142 & 0.108 & Positive & Fail \\
\bottomrule
\end{tabular}
\end{table}

The positive direction (higher BCS $\to$ higher \fone{}) is opposite to the confusability hypothesis and consistent with centrality confounding (Section~\ref{sec:finding3}).

\paragraph{DocRED.} Within-count BCS--\fone{} correlations are similarly non-significant: count 8 ($r{=}0.032$, $p{=}0.62$), count 16 ($r{=}{-}0.058$, $p{=}0.48$), pooled ($r{=}{-}0.015$, $p{=}0.84$).
The per-type BCS--\fone{} correlation is also null ($\rho{=}{-}0.04$, $p{=}0.72$).

\subsection{2$\times$2 Factorial Design: Count $\times$ Confusability Interaction}
\label{sec:factorial}

We test the interaction between count and confusability using a $2{\times}2$ factorial design on Few-NERD: \{low count$=$6, high count$=$22\} $\times$ \{low BCS, high BCS\}, with 12 batches per cell ($n{=}48$ total).

\begin{table}[t]
\centering
\caption{$2{\times}2$ cell means (\fone{}, text match) on Few-NERD. Both factors show directional effects, but the interaction is not significant.}
\label{tab:factorial_means}
\small
\begin{tabular}{lccc}
\toprule
 & \textbf{Low BCS} & \textbf{High BCS} & \textbf{BCS Effect} \\
\midrule
Low Count (6)  & 0.348 & 0.312 & $\Delta{=}{-}0.036$ \\
High Count (22) & 0.298 & 0.255 & $\Delta{=}{-}0.043$ \\
\midrule
Count Effect & $\Delta{=}{-}0.050$ & $\Delta{=}{-}0.057$ & \\
\bottomrule
\end{tabular}
\end{table}

\begin{table}[t]
\centering
\caption{Two-way ANOVA results. Count is significant; BCS is marginal; the interaction is not significant.}
\label{tab:anova}
\small
\begin{tabular}{lcccc}
\toprule
\textbf{Source} & \textbf{$F$} & \textbf{$p$} & \textbf{$\eta^2$} & \textbf{Verdict} \\
\midrule
Count       & 7.82 & 0.008 & 0.14  & Significant \\
BCS         & 2.93 & 0.094 & 0.06  & Marginal \\
Count$\times$BCS & 0.31 & 0.582 & 0.006 & NS \\
\bottomrule
\end{tabular}
\end{table}

The count main effect is significant ($F{=}7.82$, $p{=}0.008$, $\eta^2{=}0.14$), while BCS is marginal ($F{=}2.93$, $p{=}0.094$, $\eta^2{=}0.06$).
The interaction is not significant ($F{=}0.31$, $p{=}0.582$, $\eta^2{=}0.006$), consistent with an additive (rather than multiplicative) combination of the two channels.
Simple effects analysis shows that BCS effects are directionally consistent at both count levels but neither reaches significance individually (low count: $\Delta{=}{-}0.036$, $p{=}0.22$; high count: $\Delta{=}{-}0.043$, $p{=}0.18$).


\subsection{Error Taxonomy: Mechanistic Evidence for Distinct Channels}
\label{sec:error_taxonomy}

We perform a systematic error classification on Few-NERD extraction outputs, categorizing false positives into four types: type substitution, boundary error, spurious extraction, and missed entity.

\begin{table}[t]
\centering
\caption{Overall error distribution. Spurious extraction dominates; type substitution (the channel most related to confusability) accounts for only 18\%.}
\label{tab:error_dist}
\small
\begin{tabular}{lcc}
\toprule
\textbf{Error Type} & \textbf{Proportion} & \textbf{Primary Driver} \\
\midrule
Spurious Extraction & 45\% & Prompt format \\
Missed Entity       & 25\% & Count (recall) \\
Type Substitution   & 18\% & BCS (confusability) \\
Boundary Error      & 12\% & Span detection \\
\bottomrule
\end{tabular}
\end{table}

\paragraph{Type substitution by confusability.}
In extreme high-BCS batches, the type substitution rate reaches $0.058$, compared to $0.012$ in extreme low-BCS batches---a $4.8{\times}$ ratio (Table~\ref{tab:type_sub_bcs}).
However, in \emph{random} batches, the ratio drops to $1.5{\times}$ ($0.022$ vs.\ $0.015$), indicating that the effect is pronounced only under deliberately extreme conditions.

\begin{table}[t]
\centering
\caption{Type substitution rate by BCS condition. The extreme ratio (4.8$\times$) substantially exceeds the random ratio (1.5$\times$).}
\label{tab:type_sub_bcs}
\small
\begin{tabular}{lccc}
\toprule
\textbf{Condition} & \textbf{High BCS} & \textbf{Low BCS} & \textbf{Ratio} \\
\midrule
Extreme Batch & 0.058 & 0.012 & 4.8$\times$ \\
Random Batch  & 0.022 & 0.015 & 1.5$\times$ \\
\bottomrule
\end{tabular}
\end{table}

\paragraph{Missed entities by count.}
The miss rate increases monotonically with count: $0.28$ (count 6) $\to$ $0.42$ (count 22) $\to$ $0.55$ (count 64).
Crucially, the miss rate does \emph{not} vary with BCS (high BCS: $0.40$; low BCS: $0.38$; ratio $1.05$), confirming that count and confusability affect different error modes.

\begin{table}[t]
\centering
\caption{Missed entity rate by count. Monotonically increasing, independent of BCS.}
\label{tab:miss_count}
\small
\begin{tabular}{ccc}
\toprule
\textbf{Count} & \textbf{Miss Rate} & \textbf{$\Delta$ vs.\ Count 6} \\
\midrule
6  & 0.28 & --- \\
22 & 0.42 & $+$0.14 \\
64 & 0.55 & $+$0.27 \\
\bottomrule
\end{tabular}
\end{table}

\paragraph{Boundary errors and hallucination.}
Neither boundary errors (high BCS: $0.041$; low BCS: $0.038$; ratio $1.08$) nor hallucination (rate $0.03$, no BCS effect) are influenced by confusability, as expected---these error modes are orthogonal to type semantics.


% ============================================================
\section{Finding 3: Per-Type Aggregation Introduces Centrality Confound}
\label{sec:finding3}
% ============================================================

\subsection{MAVEN: Positive Per-Type Correlation}

On MAVEN, per-type mean BCS correlates \emph{positively} with per-type \fone{} under both embeddings (Table~\ref{tab:maven_pertype}).
This is the \emph{opposite} of what the confusability hypothesis predicts.

\begin{table}[t]
\centering
\caption{MAVEN pairwise and per-type confusability analysis. Per-type correlations are positive (confounded); pairwise MiniLM is significant.}
\label{tab:maven_pertype}
\small
\begin{tabular}{lcc}
\toprule
\textbf{Embedding} & \textbf{Pairwise $r$ ($p$)} & \textbf{Per-type $\rho$ ($p$)} \\
\midrule
Qwen3  & $-$0.009 (0.38) & $+$0.195 (0.018) \\
MiniLM & $+$0.115 ($2.1{\times}10^{-42}$) & $+$0.212 (0.009) \\
\bottomrule
\end{tabular}
\end{table}

The mechanism is \emph{centrality confounding}: in MAVEN's semantically homogeneous schema, types with high mean BCS are central, frequent types (e.g., \textsc{Attack}, \textsc{Movement}) that the LLM extracts well because they are well-represented in training data and have clear definitions.
Mean BCS conflates confusability with type centrality; per-type aggregation cannot separate these effects.

Crucially, the pairwise analysis tells a different story.
MiniLM pairwise SCS significantly predicts confusion ($r{=}0.115$, $p{<}10^{-42}$), demonstrating that pairwise analysis avoids the centrality confound.
The Qwen3 pairwise signal is null ($r{=}{-}0.009$), highlighting that embedding choice interacts with schema structure.

The confusion rate analysis further supports this: MAVEN's extreme confusion ratios ($5.8{\times}$ at count 6, $4.2{\times}$ at count 12) are lower than Few-NERD's ($10{-}28{\times}$), while random ratios ($1.4{\times}$, $1.3{\times}$) are modest.

\subsection{DocRED: Null Result as Boundary Condition}

DocRED yields complete null results across all confusability analyses (Section~\ref{sec:within_count}).
The pairwise signal is non-significant for both Qwen3 ($r{=}0.023$, $p{=}0.14$) and MiniLM ($r{=}0.031$, $p{=}0.08$).
The per-type BCS--\fone{} correlation is null ($\rho{=}{-}0.04$, $p{=}0.72$).

\begin{table}[t]
\centering
\caption{DocRED pairwise confusability analysis. No significant signal detected with either embedding.}
\label{tab:docred_pairwise}
\small
\begin{tabular}{lcccc}
\toprule
\textbf{Embedding} & \textbf{Pearson $r$} & \textbf{$p$} & \textbf{$n$ pairs} & \textbf{Verdict} \\
\midrule
Qwen3  & 0.023 & 0.14 & 4,560 & NS \\
MiniLM & 0.031 & 0.08 & 4,560 & NS \\
\bottomrule
\end{tabular}
\end{table}

This null result establishes a boundary condition: confusability effects require that type definitions occupy a significant fraction of the prompt.
In document-level RE, the document content dominates the prompt; type definitions are a small fraction, and performance bottlenecks lie in cross-sentence reasoning and entity pair combinatorics rather than type confusability.


% ============================================================
\section{Finding 4: Confusability-Minimizing Decomposition Fails}
\label{sec:finding4}
% ============================================================

Motivated by the confusability signal in Section~\ref{sec:finding1}, we tested whether separating confusable types into different batches (Schema DAG Decomposition, SDGD) could improve extraction.
We report a comprehensive negative result: all decomposition strategies perform worse than single-pass extraction.
Full details and per-batch breakdowns are in Appendix~\ref{app:sdgd}.

\subsection{Overview of Methods}

\begin{itemize}[leftmargin=1.5em,itemsep=2pt]
\item \textbf{Single-pass}: All types in one prompt.
\item \textbf{SDGD-zero}: Graph coloring on the SCS-thresholded pairwise similarity graph (no calibration data).
\item \textbf{SDGD-cal}: Graph coloring on an empirical confusion matrix from 8--10 calibration samples.
\item \textbf{Random ($k$)}: Random batching with the same number of batches as SDGD.
\item \textbf{Per-type}: One type per batch (UniNER-style).
\end{itemize}

\subsection{Results}

Table~\ref{tab:sdgd_main} presents the full comparison.

\begin{table}[t]
\centering
\caption{Full method comparison. All decomposition methods are inferior to single-pass. Ordering: single-pass $>$ SDGD-cal $\gtrsim$ random $\gtrsim$ SDGD-zero $>$ per-type. Methods are sorted by Few-NERD \fone{}.}
\label{tab:sdgd_main}
\small
\begin{tabular}{lccccc}
\toprule
\textbf{Method} & \multicolumn{3}{c}{\textbf{Few-NERD}} & \multicolumn{2}{c}{\textbf{MAVEN}} \\
\cmidrule(lr){2-4} \cmidrule(lr){5-6}
 & \fone{} & Prec. & Recall & Rel.\ \fone{} & API \\
\midrule
Single-pass & 0.414 & 0.648 & 0.304 & 0.241 & 1$\times$ \\
SDGD-cal    & 0.381 & 0.342 & 0.430 & 0.188 & 7--12$\times$ \\
Random      & 0.374 & 0.335 & 0.425 & 0.180 & 8--12$\times$ \\
SDGD-zero   & 0.367 & 0.318 & 0.434 & 0.168 & 8--15$\times$ \\
Per-type    & 0.330 & 0.225 & 0.620 & 0.108 & 66--168$\times$ \\
\bottomrule
\end{tabular}
\end{table}

The consistent pattern across both tasks: decomposition increases recall (each batch focuses on fewer types, enabling more aggressive extraction) but devastates precision (without the full type set, the model cannot disambiguate, leading to massive over-extraction).
The precision--recall tradeoff is unfavorable: precision losses far exceed recall gains.

\paragraph{SDGD vs.\ random is not significant.}
SDGD-zero does not outperform random batching on either task (Few-NERD: $\Delta{=}{-}0.7$pp, $p{=}0.68$; MAVEN: $\Delta{=}{-}1.2$pp, $p{=}0.42$).
Confusability-aware grouping provides no value over random partitioning.

\paragraph{Per-type splitting is worst.}
The extreme decomposition---one type per batch---yields the lowest \fone{} (Few-NERD: $0.330$; MAVEN: $0.108$) at the highest cost ($66{-}168{\times}$ API calls).
This confirms that decomposition granularity is inversely related to performance.

\paragraph{Root cause.}
Decomposition removes \emph{disambiguation context}: when confusable types are separated, each batch sees only one side of a confusable pair, and the model assigns the available type to any plausible mention.
Upon merging batches, the same span may receive conflicting type assignments ($8\%$ cross-batch span overlap), and the overall prediction count far exceeds gold annotations (pred/gold ratio $1.36$ vs.\ single-pass $0.47$).


% ============================================================
\section{Discussion}
\label{sec:discussion}
% ============================================================

\paragraph{Schema complexity is multi-dimensional but context-dependent.}
Our results establish that count and confusability are distinguishable dimensions of schema complexity, but their effects are contingent on task structure, model architecture, and embedding choice.
On Few-NERD (cross-domain NER), count dominates; on MAVEN (same-domain ED), more types \emph{help}; on DocRED (document-level RE), neither dimension matters.

\paragraph{Confusability effects are real but not universal.}
The pairwise confusability signal is significant for Qwen3 on Few-NERD and MiniLM on MAVEN, but absent for LLaMA-3.3-70B and for Qwen3 embeddings on MAVEN.
This model dependence suggests that confusability sensitivity is a property of the model's internal representation, not solely of the schema itself.
The $4{\times}$ effect size difference between Qwen3 ($r{=}0.156$) and LLaMA ($r{=}0.042$) is itself a finding that may explain prior conflicting reports about schema complexity effects.

\paragraph{Per-type aggregation is unreliable.}
The centrality confound on MAVEN demonstrates that mean BCS (per-type averaged confusability) is not a valid measure in semantically homogeneous schemas.
Future work on schema complexity should use pairwise analysis to avoid this confound, or use empirical confusion matrices that bypass embedding proxy limitations.

\paragraph{Decomposition harms LLM IE.}
The most counter-intuitive finding is that confusability-minimizing decomposition \emph{hurts} performance.
LLMs use type co-presence for comparative reasoning: seeing confusable types together enables the model to reason about their distinctions.
Removing this context forces the model into uninformed guessing.
This insight has direct implications for decomposition-based approaches: they should be designed to \emph{preserve} disambiguation context rather than minimize confusability.

\paragraph{Implications for future work.}
(1)~Confusability-based methods should be evaluated per-model, not assumed universal.
(2)~Schema decomposition for LLM IE needs fundamentally different principles---perhaps preserving confusable types together while reducing only the non-informative count overhead.
(3)~Embedding-based confusability proxies need task-specific calibration, as no single embedding works universally.

\paragraph{Limitations.}
Our study uses two LLM families; more models are needed to characterize the distribution of confusability sensitivity.
The $2{\times}2$ factorial design has limited statistical power ($n{=}48$); larger-scale factorial experiments could detect smaller interaction effects.
DocRED results may partly reflect the difficulty of zero-shot document-level RE rather than a true absence of confusability effects.


% ============================================================
\section{Conclusion}
\label{sec:conclusion}
% ============================================================

We presented the first controlled empirical study isolating type count and type confusability as dimensions of schema complexity in training-free LLM information extraction.
Our experiments across three IE tasks, two LLMs, and two embedding spaces reveal that:
(1)~pairwise type confusability is a valid predictor of type confusion errors, but only for specific model$\times$embedding combinations;
(2)~count and confusability degrade performance through distinguishable error channels;
(3)~per-type confusability aggregation introduces centrality confounds in homogeneous schemas; and
(4)~confusability-minimizing decomposition is counter-productive because LLMs leverage type co-presence for disambiguation.
These findings reframe schema complexity as a multi-dimensional, model-dependent phenomenon and provide methodological guidance for future research on LLM IE with complex schemas.


% ============================================================
\appendix
\section{SDGD Decomposition: Detailed Results}
\label{app:sdgd}
% ============================================================

This appendix provides detailed results for all decomposition experiments, which are reported as negative results in the main paper (Section~\ref{sec:finding4}).

\subsection{SDGD-zero on Few-NERD}
\label{app:sdgd_zero_fewnerd}

Graph coloring on the Qwen3 SCS-thresholded graph ($\tau{=}0.847$) produces 8 batches (avg.\ size $8.25$, range $4{-}12$).

\begin{table}[h]
\centering
\caption{SDGD-zero vs.\ single-pass on Few-NERD (66 types).}
\label{tab:sdgd_zero_fewnerd}
\small
\begin{tabular}{lccccc}
\toprule
\textbf{Method} & \textbf{\fone{}} & \textbf{Precision} & \textbf{Recall} & \textbf{Pred/Gold} & \textbf{API} \\
\midrule
Single-pass & 0.414 & 0.648 & 0.304 & 0.47 & 1$\times$ \\
SDGD-zero   & 0.367 & 0.318 & 0.434 & 1.36 & 8$\times$ \\
$\Delta$    & $-$4.7pp & $-$33.0pp & $+$13.0pp & $+$0.89 & $+$7$\times$ \\
\bottomrule
\end{tabular}
\end{table}

\begin{table}[h]
\centering
\caption{SDGD-zero per-batch breakdown (top 4 batches, Few-NERD).}
\label{tab:sdgd_zero_perbatch}
\small
\begin{tabular}{cccccc}
\toprule
\textbf{Batch} & \textbf{Types} & \textbf{Avg BCS} & \textbf{\fone{}} & \textbf{Recall} & \textbf{Precision} \\
\midrule
0 & 12 & 0.812 & 0.42 & 0.48 & 0.38 \\
1 & 10 & 0.795 & 0.38 & 0.45 & 0.33 \\
2 & 9  & 0.780 & 0.35 & 0.42 & 0.30 \\
3 & 8  & 0.763 & 0.32 & 0.40 & 0.27 \\
\bottomrule
\end{tabular}
\end{table}

Graph statistics: 142 edges, edge density $0.066$, chromatic number 8, max intra-batch BCS $0.842$.
Cross-batch span overlap: $8\%$ of predictions appear in multiple batches.
Total predictions: $4{,}850$ (after dedup: $4{,}120$), averaging $8.2$ predictions per sample vs.\ $5.9$ gold entities.

\subsection{SDGD-zero on MAVEN}
\label{app:sdgd_zero_maven}

On MAVEN (168 types), adaptive threshold $\tau{=}0.319$ produces 15 batches (avg.\ size $11.2$, range $5{-}18$), with edge density $0.132$ and $1{,}850$ edges.

\begin{table}[h]
\centering
\caption{SDGD-zero vs.\ single-pass on MAVEN (168 types, relaxed matching).}
\label{tab:sdgd_zero_maven}
\small
\begin{tabular}{lcccc}
\toprule
\textbf{Method} & \textbf{Relaxed \fone{}} & \textbf{Exact \fone{}} & \textbf{API} & \textbf{Wall-Clock} \\
\midrule
Single-pass & 0.241 & 0.018 & 1$\times$ & 1$\times$ \\
SDGD-zero   & 0.168 & 0.012 & 15$\times$ & 12.8$\times$ \\
$\Delta$    & $-$7.3pp & $-$0.6pp & $+$14$\times$ & $+$11.8$\times$ \\
\bottomrule
\end{tabular}
\end{table}

The larger performance gap on MAVEN ($-$7.3pp vs.\ $-$4.7pp on Few-NERD) is consistent with MAVEN's positive count curve: decomposition directly counteracts the benefit of full ontology visibility.

\subsection{SDGD-cal on Few-NERD}
\label{app:sdgd_cal_fewnerd}

SDGD-cal uses 8 calibration samples to build an empirical confusion matrix, yielding 85 nonzero confusion pairs.
The empirical confusion matrix correlates weakly with embedding SCS ($\rho{=}0.125$, $p{=}0.031$).
Graph coloring produces 7 batches (avg.\ size $9.4$).

\begin{table}[h]
\centering
\caption{SDGD-cal vs.\ baselines on Few-NERD.}
\label{tab:sdgd_cal_fewnerd}
\small
\begin{tabular}{lccccc}
\toprule
\textbf{Method} & \textbf{\fone{}} & \textbf{$\Delta$ vs.\ SP} & \textbf{Precision} & \textbf{Recall} & \textbf{API} \\
\midrule
Single-pass & 0.414 & --- & 0.648 & 0.304 & 1$\times$ \\
Random ($k{=}8$) & 0.374 & $-$4.0 & 0.335 & 0.425 & 8$\times$ \\
SDGD-zero   & 0.367 & $-$4.7 & 0.318 & 0.434 & 8$\times$ \\
SDGD-cal    & 0.381 & $-$3.3 & 0.342 & 0.430 & 7$\times$ $+$ 8 cal. \\
\bottomrule
\end{tabular}
\end{table}

SDGD-cal marginally outperforms SDGD-zero ($+$1.4pp) and random ($+$0.7pp, $p{=}0.58$), but the improvement over random is not significant.
Even with empirical confusion data, the decomposition penalty persists.

\subsection{SDGD-cal on MAVEN}
\label{app:sdgd_cal_maven}

On MAVEN, 10 calibration samples produce 210 nonzero confusion pairs.
Graph coloring yields 12 batches (avg.\ size $14.0$).

\begin{table}[h]
\centering
\caption{Full comparison on MAVEN (168 types, relaxed matching).}
\label{tab:sdgd_cal_maven}
\small
\begin{tabular}{lccc}
\toprule
\textbf{Method} & \textbf{Relaxed \fone{}} & \textbf{$\Delta$ vs.\ SP} & \textbf{API} \\
\midrule
Single-pass  & 0.241 & --- & 1$\times$ \\
Random ($k{=}12$) & 0.180 & $-$6.1 & 12$\times$ \\
SDGD-zero    & 0.168 & $-$7.3 & 15$\times$ \\
SDGD-cal     & 0.188 & $-$5.3 & 12$\times$ $+$ 10 cal. \\
\bottomrule
\end{tabular}
\end{table}

SDGD-cal outperforms SDGD-zero ($+$2.0pp) but not random ($+$0.8pp, $p{=}0.62$).
The same pattern as Few-NERD: calibration helps marginally, but decomposition remains inferior to single-pass.

\subsection{Random Batching Ablation}
\label{app:random}

To isolate the contribution of confusability-aware grouping from the batch count reduction effect, we compare SDGD-zero against random batching with the same number of batches.

\begin{table}[h]
\centering
\caption{Random vs.\ SDGD-zero on Few-NERD ($k{=}8$, 10 random seeds).}
\label{tab:random_fewnerd}
\small
\begin{tabular}{lcccc}
\toprule
\textbf{Method} & \textbf{Mean \fone{}} & \textbf{Std} & \textbf{Range} & \textbf{vs.\ SP} \\
\midrule
Random$\times$10 & 0.374 & 0.018 & [0.348, 0.402] & $-$4.0pp \\
SDGD-zero & 0.367 & --- & --- & $-$4.7pp \\
$\Delta$ (SDGD$-$Random) & \multicolumn{4}{c}{$-$0.7pp, $p{=}0.68$ (NS)} \\
\bottomrule
\end{tabular}
\end{table}

\begin{table}[h]
\centering
\caption{Random vs.\ SDGD-zero on MAVEN.}
\label{tab:random_maven}
\small
\begin{tabular}{lccc}
\toprule
\textbf{Method} & \textbf{Relaxed \fone{}} & \textbf{vs.\ SP} & \textbf{$p$ (paired)} \\
\midrule
Random$\times$5 & 0.180 & $-$6.1pp & --- \\
SDGD-zero & 0.168 & $-$7.3pp & --- \\
$\Delta$ & \multicolumn{3}{c}{$-$1.2pp, $p{=}0.42$ (NS)} \\
\bottomrule
\end{tabular}
\end{table}

On both tasks, SDGD-zero does not significantly outperform random batching, confirming that confusability-aware grouping provides no measurable benefit.
SDGD-zero may actually perform slightly \emph{worse} than random, because deliberately separating confusable types removes same-cluster disambiguation anchors that random batching occasionally preserves.

\subsection{Per-Type Splitting Ablation}
\label{app:pertype}

Per-type splitting (one type per batch) represents the extreme of decomposition---maximum confusability elimination but also maximum loss of disambiguation context.

\begin{table}[h]
\centering
\caption{Per-type splitting vs.\ all methods on Few-NERD.}
\label{tab:pertype_fewnerd}
\small
\begin{tabular}{lccccc}
\toprule
\textbf{Method} & \textbf{\fone{}} & \textbf{Precision} & \textbf{Recall} & \textbf{API} & \textbf{Wall-Clock} \\
\midrule
Single-pass & 0.414 & 0.648 & 0.304 & 1$\times$ & 1$\times$ \\
SDGD-cal    & 0.381 & 0.342 & 0.430 & 7$\times$ & 5.5$\times$ \\
Random ($k{=}8$) & 0.374 & 0.335 & 0.425 & 8$\times$ & 6.5$\times$ \\
SDGD-zero   & 0.367 & 0.318 & 0.434 & 8$\times$ & 6.5$\times$ \\
Per-type    & 0.330 & 0.225 & 0.620 & 66$\times$ & 48$\times$ \\
\bottomrule
\end{tabular}
\end{table}

\begin{table}[h]
\centering
\caption{Per-type splitting vs.\ all methods on MAVEN (relaxed matching).}
\label{tab:pertype_maven}
\small
\begin{tabular}{lccc}
\toprule
\textbf{Method} & \textbf{Relaxed \fone{}} & \textbf{API} & \textbf{Wall-Clock} \\
\midrule
Single-pass & 0.241 & 1$\times$ & 1$\times$ \\
SDGD-cal    & 0.188 & 12$\times$ & 10$\times$ \\
Random      & 0.180 & 12$\times$ & 10$\times$ \\
SDGD-zero   & 0.168 & 15$\times$ & 12.8$\times$ \\
Per-type    & 0.108 & 168$\times$ & 135$\times$ \\
\bottomrule
\end{tabular}
\end{table}

The method ordering is consistent across both tasks: single-pass $>$ SDGD-cal $\gtrsim$ random $\gtrsim$ SDGD-zero $>$ per-type.
Decomposition granularity is inversely correlated with \fone{}: larger batches (more type context) consistently outperform smaller ones.
Per-type splitting achieves the highest recall (0.620 on Few-NERD) through aggressive per-type extraction, but precision collapses (0.225) as each batch drastically over-extracts.

% ============================================================
% References
% ============================================================

\bibliographystyle{plainnat}

\begin{thebibliography}{10}

\bibitem{universal_ner_2024}
W.~Zhou, S.~Zhang, Y.~Gu, M.~Chen, and H.~Peng.
\newblock Universal{NER}: Targeted distillation from large language models for open named entity recognition.
\newblock In \emph{Proc.\ ICLR}, 2024.

\bibitem{gollie_2024}
O.~Sainz, I.~Garc{\'i}a-Ferrero, R.~Agerri, O.~L{\'o}pez~de~Lacalle, G.~Rigau, and E.~Agirre.
\newblock Go{LLIE}: Annotation guidelines improve zero-shot information extraction.
\newblock In \emph{Proc.\ ICLR}, 2024.

\bibitem{llmpee_2024}
G.~Zaratiana, N.~Tomeh, P.~Holat, and T.~Charnois.
\newblock {LLM-PEE}: Prompt-based extraction with embedding similarity for entity extraction using large language models.
\newblock \emph{arXiv preprint arXiv:2312.15805}, 2024.

\bibitem{few_nerd_2021}
N.~Ding, G.~Xu, Y.~Chen, X.~Wang, X.~Han, P.~Xie, H.~Zheng, and Z.~Liu.
\newblock Few-{NERD}: A few-shot named entity recognition dataset.
\newblock In \emph{Proc.\ ACL-IJCNLP}, 2021.

\bibitem{maven_2020}
X.~Wang, Z.~Wang, X.~Han, W.~Jiang, R.~Han, Z.~Liu, J.~Li, P.~Li, Y.~Lin, and J.~Zhou.
\newblock {MAVEN}: A massive general domain event detection dataset.
\newblock In \emph{Proc.\ EMNLP}, 2020.

\end{thebibliography}

\end{document}
